% !TEX root = ../main.tex
% SY02-29 §8 Kottwitz 猜想（原文 p.36–p.42）
% 译手：kottwitz
% 说明：本块只承载原文 §8；p.42 顶部即为 §9（归 s11），故 \srcpage{42} 置于文件末尾。

\section{Kottwitz 猜想：基本 RZ 空间的上同调（[22]）}
\srcpage{36}

\subsection{基本例子：Lubin--Tate 空间与 Drinfeld 空间}

设 $G=\mathrm{GL}_n$ 定义在 $p$-进域 $E$ 上。并令 $\mu(z)=\mathrm{diag}(z,1,\dots,1)$。固定
\[
b=\begin{pmatrix}
0&\cdots&0&\pi\\
1&&&0\\
&\ddots&&\vdots\\
&&1&0
\end{pmatrix}^{-1}
\]
使得 $\mathscr{E}_b=\oo(\frac1n)$，且 $J_b(E)=D^\times$，其中 $D$ 是 $E$ 上不变量为 $1/n$ 的除代数。我们用 $\breve{E}$ 记 $E$ 的最大非分歧扩域的完备化。\yizhuw{$b$ 取此形状，是因为它对应 $B(\mathrm{GL}_n)$ 中斜率 $1/n$ 的基本元素：右上角的 $\pi$ 记录平移，$\mathscr{E}_b=\oo(\frac1n)$ 是其稳定向量丛，$J_b$ 是它的 $\sigma$-中心化子，在 $\mathrm{GL}_n$ 情形即 $D^\times$。}

有笛卡尔图表
\begin{equation*}
\begin{tikzcd}
{[\PP^{n-1,\diamond}_E/\mathrm{GL}_n(E)]} \arrow[r] \arrow[d,"i"']
& {[\Spa(E)/\mathrm{GL}_n(E)]} \arrow[d,"{(x_1,\id)}"] \\
\mathrm{Hecke}^{\mu} \arrow[r,"\overrightarrow{h}"]
& \Bun_G\times\Spa(E)^{\diamond}.
\end{tikzcd}
\end{equation*}

这里 $x_1$ 由平凡向量丛给出。事实上，给自己一个如下形式的平凡向量丛的修正
\begin{equation}\label{eq:8.9}\tag{9}
0\longrightarrow\oo^n\longrightarrow\mathscr{E}\longrightarrow i_*\mathscr{F}\longrightarrow 0,
\end{equation}
其中 $\mathscr{F}$ 在 $R^\sharp$ 上局部自由、秩为 $1$，就等于给自己一个元素 $\PP^{n-1}_E(R^\sharp)$（如前，$R^\sharp$ 是 $R$ 的一个 untilt，$i$ 是对应的 Cartier 除子 $\Spec(R^\sharp)\hookrightarrow X_R^{sch}$）。复合态射
\[
f:[\PP^{n-1,\diamond}_E/\mathrm{GL}_n(E)]\longrightarrow\mathit{Hecke}^{\mu}\xrightarrow{\overleftarrow{h}}\Bun_G
\]
给出前述修正 (9) 中的向量丛 $\mathscr{E}$。几何地逐纤维看，$\mathscr{E}$ 的同构类为 $\oo^i\oplus\oo(\frac{1}{n-i})$，其中整数 $i\in\{0,\dots,n-1\}$。$\PP^{n-1}$ 上相应的分层是 $\Bun_G$ 的 HN 分层沿 $f$ 的拉回。开层对应\srcpageinl{37}的是这样的轨迹：在每条几何纤维上 $\mathscr{E}$ 同构于 $\oo(\frac1n)$。它是半稳定轨迹沿 $f$ 的拉回，并且与 Drinfeld 空间 $\Omega$ 重合；更确切地说，这就是开子集
\[
[\Omega^{\diamond}/\mathrm{GL}_n(E)]\subset[\PP^{n-1,\diamond}_E/\mathrm{GL}_n(E)].
\]
有相应的态射
\[
f|_{[\Omega^{\diamond}/\mathrm{GL}_n(E)]}:[\Omega^{\diamond}_{\breve{E}}/\mathrm{GL}_n(E)]\longrightarrow\Bun^{ss}_{G,\FF_q}\xleftarrow[x_b]{\sim}[\Spa(\ol{\FF}_q)/D^\times].
\]
沿这个态射拉回 $[\Spa(\FF_q)/D^\times]$ 上的万有 $D^\times$-挠子，就得到 Drinfeld 覆盖 $\Omega_{\breve{E}}$ 的 Rapoport--Zink 版本。在这座塔的顶端，前述把 $\oo^n$ 修正为 $\oo(\frac1n)$ 的修正 (9)，正是无穷水平 Lubin--Tate 塔上万有 $p$-可除群所对应的 Hodge--Tate 正合列。\yizhuw{Lubin--Tate 塔与 Drinfeld 塔在无穷水平互为对偶（$p$-可除群与形式 $\oo_D$-模互换 $G$ 与 $D^\times$），所以“Drinfeld 覆盖的 Rapoport--Zink 版本”是把同一族挠子放到 RZ 一侧来读。}

设 $\varphi:W_E\longrightarrow\mathrm{GL}_n(\ol{\QQ}_\ell)$ 是离散 Langlands 参数（在此情形即不可分解表示）。设 $\pi$ 是局部 Langlands 对应下 $\mathrm{GL}_n(E)$ 的相应平方可积表示，$\rho$ 是 Jacquet--Langlands 对应下 $D^\times$ 的相应表示，$\mathrm{JL}(\rho)=\pi$。按猜想，
\[
\begin{aligned}
x_1^*\mathscr{F}_{\varphi}&=\pi\\
x_b^*\mathscr{F}_{\varphi}&=\rho.
\end{aligned}
\]
把真基变换用于 Hecke 性质，得到
\[
R\Gamma_c\big(\PP^{n-1}_{\mathbb{C}_p}/\mathrm{GL}_n(E),i^*\overleftarrow{h}^*\mathscr{F}_{\varphi}\big)\big(\tfrac{n-1}{2}\big)[n-1]=\pi\otimes\varphi.
\]
这是说，$i^*\mathscr{F}_{\varphi}$ 的光滑等变上同调复形等于 $\pi\otimes\varphi(\frac{1-n}{2})$，并集中在中间次数。

现在再设 $\varphi$ 是尖的，即表示 $\varphi$ 不可约。猜想中的尖性条件于是说：若 $j:\Omega\hookrightarrow\PP^{n-1}$，则
\[
i^*\overleftarrow{h}^*\mathscr{F}_{\varphi}=j_!j^*i^*\overleftarrow{h}^*\mathscr{F}_{\varphi}.
\]
但此时，若 $\mathcal{M}_\infty\to\Omega_{\breve{E}}$ 是无穷水平的 Drinfeld 塔，则
\[
j^*i^*\overleftarrow{h}^*\mathscr{F}_{\varphi}=\mathcal{M}_\infty\times_{D^\times}\rho.
\]
猜想于是预言
\[
R\Gamma_c(\mathcal{M}_\infty\otimes\mathbb{C}_p,\ol{\QQ}_\ell)\otimes_{D^\times}\rho=\pi\otimes\varphi\big(\tfrac{1-n}{2}\big)[1-n],
\]
这正是局部 Langlands/Jacquet--Langlands 对应在 Drinfeld 塔上同调中众所周知的实现。

\begin{remark}\label{rem:8.1}
人们会注意到：对离散而非尖的 $\varphi$，猜想预言 $[\PP^{n-1}/\mathrm{GL}_n(E)]$ 上存在一个反常层，其上同调恰在中间次数实现局部 Langlands 对应。这与通常用 Drinfeld 塔上同调实现局部 Langlands 对应的方式不同（一般而言，对非超尖参数，相应表示不会出现在中间次数）。这个反常层在 $\Omega$ 上的限制应当是 $\mathcal{M}_\infty\times_{D^\times}\rho$。由此可以把 $[\PP^{n-1}/\mathrm{GL}_n(E)]$ 上以该反常层为系数的上同调，看作 $\Omega$ 上以 $\mathcal{M}_\infty\times_{D^\times}\rho$ 为系数的相交上同调。
\end{remark}

\subsection{一般情形}

\srcpage{38}

\subsubsection*{8.2.1 局部 shtuka 模空间}

我们将看到，这个猜想蕴含 Kottwitz 猜想，后者描述基本 Rapoport--Zink 空间上同调的超尖部分。以下取 $E=\QQ_p$，从而 $G/\QQ_p$。

固定 $b\in G(L)$ 与 $\mu\in X_*(T)^+/\Gamma$。我们打算考察 $\mu$ 的某个 Galois 轨道上的 Rapoport--Zink 空间的非交并。好处是无需引入任何反射域，这个非交并便定义在 $\QQ_p$ 上（这简化了记号）。当 $L$ 是我们模空间的基域时，记 $L=\breve{\QQ}_p$。

来看下面的图表
\begin{equation*}
\begin{tikzcd}
\Spa(\ol{\FF}_p) \arrow[dr,"x_b"']
& & \mathrm{Hecke}^{\le\mu}\otimes\ol{\FF}_p \arrow[dl,"\overleftarrow{h}"] \arrow[dr,"\overrightarrow{h}"]
& & \Spa(\breve{\QQ}_p)^{\diamond} \arrow[dl,"{(x_1,\id)}"] \\
& \Bun_{G,\ol{\FF}_p} & & \Bun_{G,\ol{\FF}_p}\times\Spa(\breve{\QQ}_p)^{\diamond} &
\end{tikzcd}
\end{equation*}
并把相伴的局部 shtuka 模空间定义为该图表的纤维积：
\begin{gather*}
\Sht(G,b,\mu):=\Spa(\ol{\FF}_p)\times_{x_b,\,\Bun_{G,\ol{\FF}_p},\,\overleftarrow{h}}\mathit{Hecke}^{\le\mu}\otimes\ol{\FF}_p\\
\times_{\overrightarrow{h},\,\Bun_{G,\ol{\FF}_p}\times\Spa(\breve{\QQ}_p)^{\diamond},\,(x_1,\id)}\Spa(\breve{\QQ}_p)^{\diamond}.
\end{gather*}
这是 $\breve{\QQ}_p^{\diamond}=\QQ_p^{\diamond}\otimes_{\FF_p}\ol{\FF}_p$ 上的菱形。它带有 $J_b(\QQ_p)$ 的作用，来自分解
\[
x_b:\Spa(\ol{\FF}_p)\longrightarrow[\Spa(\ol{\FF}_p)/J_b(\QQ_p)]\longrightarrow\Bun_{G,\ol{\FF}_p},
\]
还带有与之交换的 $G(\QQ_p)$ 的作用，来自分解
\[
x_1:\Spa(\ol{\FF}_p)\longrightarrow[\Spa(\ol{\FF}_p)/G(\QQ_p)]\longrightarrow\Bun_{G,\ol{\FF}_p}.
\]

带 $G(\QQ_p)\times J_b(\QQ_p)$ 作用的菱形 $\Sht(G,b,\mu)$ 具有从 $\breve{\QQ}_p^{\diamond}$ 到 $\QQ_p$ 的\emph{Weil 下降数据}。事实上，$x_1$ 定义在 $\FF_p$ 上。而 $x_b$ 不定义在 $\FF_p$ 上，但它的同构类定义在 $\FF_p$ 上。更精确地，$\Frob^*x_b=x_{b^{\sigma}}$，且有图表
\begin{equation*}
\begin{tikzcd}
\Spa(\ol{\FF}_p) \arrow[rr,bend left=20,"x_{b^{\sigma}}"] \arrow[rr,bend right=20,"x_b"'] & {\Downarrow} & \Bun_{G,\ol{\FF}_p}
\end{tikzcd}
\end{equation*}
其中竖直箭头由 $\sigma$-共轭 $b^{\sigma}=b^{-1}.b.b^{\sigma}$ 给出。\yizhu{原文如此：等号两端同为 $b^{\sigma}$，字面自指。$\sigma$-共轭的本义是 $b\mapsto b^{-1}b\,b^{\sigma}$（参 §2.2 的定义）。译文照排，未擅改。}{A：原文 p.38 与 §2.2 的 $\sigma$-共轭定义；已登记疑点清单}

当然，要使 $\Sht(G,b,\mu)$ 非空，必须假设 $[b]\in B(G,\mu)$。有两个周期态射
\begin{equation*}
\begin{tikzcd}
& \Sht(G,b,\mu) \arrow[dl,"\pi_{dR}"'] \arrow[dr,"\pi_{HT}"] & \\
\Gr^{\le-\mu}\otimes_{\QQ_p^{\diamond}}\breve{\QQ}_p^{\diamond} & & \Gr^{\le\mu}
\end{tikzcd}
\end{equation*}
分别由 $\overleftarrow{h}$ 与 $\overrightarrow{h}$ 诱导。态射 $\pi_{dR}$（相应地 $\pi_{HT}$）是 $G(\QQ_p)$-不变的（相应地 $J_b(\QQ_p)$-不变的），并且与 $J_b(\QQ_p)$（相应地 $G(\QQ_p)$）的作用交换。若假设 $b$ 是基本的且 decent，并沿 $b\sigma$ 扭曲 $\Gr^{\le-\mu}\otimes\breve{\QQ}_p^{\diamond}$，则左端可以定义在 $\QQ_p$ 上（在极小（minuscule）情形，这给出 $G/P_\mu$ 的扭形式，例如 Severi--Brauer 簇）。这不过是核对与前述下降数据相容性的一道练习。\yizhuw{“decent”是原文沿用的技术形容词，指 $b$ 的 $\sigma$-共轭类具有某种整性（可由整值代表元实现）；原文未给定义，译文保留原文。}

\srcpage{39}

下面经由所谓的可容许轨迹，给出 $\Sht(G,b,\mu)$ 的两步构造。考虑图表
\begin{equation*}
\begin{tikzcd}
\Gr^{\le-\mu}\otimes\breve{\QQ}_p^{\diamond} \arrow[r] \arrow[d]
& \mathrm{Hecke}^{\le\mu}\otimes\ol{\FF}_p \arrow[d,"\overleftarrow{h}"'] \arrow[dr,"\overrightarrow{h}"] & \\
\Spa(\ol{\FF}_p) \arrow[r,"x_b"']
& \Bun_{G,\ol{\FF}_p} & \Bun_{G,\ol{\FF}_p}\times\Spa(\breve{\QQ}_p)^{\diamond}
\end{tikzcd}
\end{equation*}
其中方块是笛卡尔的。这定义了一个态射
\[
f:\Gr^{\le-\mu}\otimes\breve{\QQ}_p^{\diamond}\longrightarrow\mathit{Hecke}^{\le\mu}\otimes\ol{\FF}_p\xrightarrow{\overrightarrow{h}}\Bun_{G,\ol{\FF}_p}\times\Spa(\breve{\QQ}_p)^{\diamond}\xrightarrow{\mathrm{proj}}\Bun_{G,\ol{\FF}_p}.
\]
定义
\[
\big(\Gr^{\le-\mu}_{\breve{\QQ}_p^{\diamond}}\big)^{ad}=f^{-1}\big(\Bun^{1,ss}_{G,\ol{\FF}_p}\big),
\]
即所谓的可容许轨迹。这是一个开子集。\yizhuw{“可容许”是 Rapoport--Zink 理论的术语：它恰是 $\pi_{dR}$ 的定义域，而 $\pi_{HT}$ 把它映到 Newton 层 $\Gr^{\le\mu,b}$ 上，$\Sht$ 由此覆盖后者。}于是有态射
\[
\big(\Gr^{\le-\mu}_{\breve{\QQ}_p^{\diamond}}\big)^{ad}\longrightarrow\Bun^{1,ss}_{G,\ol{\FF}_p}=[\Spa(\ol{\FF}_p)/\ul{G(\QQ_p)}].
\]
这意味着可容许轨迹 $\big(\Gr^{\le-\mu}_{\breve{\QQ}_p^{\diamond}}\big)^{ad}$ 上有一个 $\ul{G(\QQ_p)}$-pro-étale 挠子。于是 $\Sht(G,b,\mu)$ 就是该 $\ul{G(\QQ_p)}$-挠子的平凡化模空间。\yizhuw{原文在此处及下文若干处给 $G(\QQ_p)$、$J_b(\QQ_p)$ 加下划线（如本句），未加说明；译文照原样保留，不据此改动任何数学内容。}

我们有以下性质：
\begin{itemize}
\item $\pi_{dR}$ 是 $\big(\Gr^{\le-\mu}_{\breve{\QQ}_p^{\diamond}}\big)^{ad}$ 上的 $\ul{G(\QQ_p)}$-挠子；
\item $\pi_{HT}$ 是 Newton 层 $\Gr^{\le\mu,b}_{\breve{\QQ}_p^{\diamond}}\subset\Gr^{\le-\mu}_{\breve{\QQ}_p^{\diamond}}$ 上的 $\widetilde{J}_b$-挠子。
\end{itemize}

现在回到猜想。设 $b$ 是基本的。我们把 Hecke 性质写成
\[
\overrightarrow{h}_!\big(\overleftarrow{h}^*\mathscr{F}_{\varphi}\otimes IC_\mu\big)=\mathscr{F}_{\varphi}\boxtimes r_\mu\circ\varphi.
\]
\yizhuw{原文在此未定义 $IC_\mu$：它是 $\Gr^{\le\mu}$ 上的相交上同调复形。这条 Hecke 性质是说，沿 $\overrightarrow{h}$ 作正向推前，$\mu$ 那部分被“算进来”。}
把真基变换用于下面的笛卡尔图表
\begin{equation*}
\begin{tikzcd}
{[\Gr^{\le\mu}\otimes\breve{\QQ}_p^{\diamond}/\ul{G(\QQ_p)}]} \arrow[r] \arrow[d,"i"']
& {[\Spa(\breve{\QQ}_p)^{\diamond}/\ul{G(\QQ_p)}]} \arrow[d,"{(x_1,\id)}"] \\
\mathrm{Hecke}^{\le\mu}\otimes\ol{\FF}_p \arrow[r,"\overrightarrow{h}"]
& \Bun_{G,\ol{\FF}_p}\times\Spa(\breve{\QQ}_p)^{\diamond}
\end{tikzcd}
\end{equation*}
沿上方水平箭头投影，就是光滑等变上同调复形。我们把该复形记作 $R\Gamma_c(\Gr^{\le\mu}_{\mathbb{C}_p^{\flat}}/G(\QQ_p),-)$，它是带 $G(\QQ_p)$ 的光滑作用、并带 $I_{\QQ_p}\subset W_{\QQ_p}$ 作用的 $\ol{\QQ}_\ell$-向量空间导出范畴中的一个对象。于是得到
\[
R\Gamma_c\big(\Gr^{\le-\mu}_{\mathbb{C}_p^{\flat}}/G(\QQ_p),(i^*\overleftarrow{h}^*\mathscr{F}_{\varphi})\otimes IC_\mu\big)=x_1^*\mathscr{F}_{\varphi}\otimes r_\mu\circ\varphi.
\]
按假设有
\[
x_1^*\mathscr{F}_{\varphi}=\bigoplus_{\rho\in\pi_0(S_{\varphi})}\rho\otimes\pi_\rho,
\]
其中 $\{\pi_\rho\}_\rho$ 是局部 Langlands 对应下与 $\varphi$ 相联系的 $G(\QQ_p)$ 表示的 L-包。\yizhuw{$S_\varphi$ 是 $\varphi$ 在 $\widehat{G}$ 中中心化子的分量群；对离散参数，$\pi_0(S_\varphi)$ 就是 L-包的指标集，所以 $x_1^*\mathscr{F}_\varphi$ 按 L-包的成员分解。}

\srcpage{40}

现在再设 $\varphi$ 是尖的、$b$ 是基本的，即 $[b]$ 是 $B(G,\mu)$ 中唯一的基本元。由于 $[b]\in B(G,\mu)$，有分解
\begin{equation*}
\begin{tikzcd}
\Gr^{\le\mu}_{\breve{\QQ}_p^{\diamond}} \arrow[r,"i"] \arrow[drr]
& \mathrm{Hecke}^{\le\mu}_{\ol{\FF}_p} \arrow[r,"\overleftarrow{h}"]
& \Bun_{G,\ol{\FF}_p} & {} \\
\Gr^{\le\mu,b}_{\breve{\QQ}_p^{\diamond}} \arrow[u,"j",hook] \arrow[drr]
& {} & \Bun^{\kappa(b)}_{G,\ol{\FF}_p} \arrow[u,hook] & {} \\
{} & {} & \Bun^{\kappa(b),ss}_{G,\ol{\FF}_p} \arrow[u,hook]
& {[\Spa(\ol{\FF}_p)/\ul{J_b(\QQ_p)}]} \arrow[l,"x_b"',"\sim" above]
\end{tikzcd}
\end{equation*}
其中各包含映射都是开子集。\yizhuw{$\kappa(b)$ 是 $b$ 的 Kottwitz 不变量，取值于 $Z(\widehat{G})^{\Gamma}$ 的特征群；它决定 $b$ 落在哪个 Newton 层，也是下文按 $\rho|_{Z(\widehat{G})^{\Gamma}}=\kappa(b)$ 筛选 L-包成员的依据。}利用尖性条件，通过 $\ul{J_b(\QQ_p)}$-挠子 $\pi_{HT}:\Sht(G,b,\mu)\to\Gr^{\le-\mu}_{\breve{\QQ}_p^{\diamond}}$，我们得到
\[
i^*\overleftarrow{h}^*\mathscr{F}_{\varphi}=j_!\Big(\Sht(G,b,\mu)\times_{J_b(\QQ_p)}x_b^*\mathscr{F}_{\varphi}\Big).
\]
利用公式
\[
x_b^*\mathscr{F}_{\varphi}=\bigoplus_{\substack{\rho\in\widehat{S_{\varphi}}\\ \rho|_{Z(\widehat{G})^{\Gamma}}=\kappa(b)}}\rho\otimes\pi_\rho
\]
我们得到以下结果（必须用特征层性质核对与 Frobenius 下降数据作用的相容性，这留给勇敢的读者……）。为得到这个公式，用到前述所有同构都与 $S_{\varphi}$ 的作用交换。

\begin{proposition}\label{prop:8.2}
猜想预言如下。取 $\mu\in X_*(T)^+/\Gamma$，设 $[b]\in B(G,\mu)$ 是唯一的基本元。设 $\varphi:W_{\QQ_p}\to{}^LG$ 是离散尖参数。给定满足 $\rho|_{Z(\widehat{G})^{\Gamma}}=\kappa(b)$ 的 $\rho\in\widehat{S_{\varphi}}$，记 $\pi_\rho$ 为与 $\varphi$ 相联系的 L-包中 $J_b(\QQ_p)$ 的相应超尖表示。考虑如下分解
\[
\mathscr{H}om(\rho,r_\mu\circ\varphi)=\bigoplus_{\rho'\in(S_{\varphi}/Z(\widehat{G})^{\Gamma})^{\wedge}}\rho'\otimes\sigma_{\rho,\rho',\mu,\varphi}
\]
作为 $S_{\varphi}\times W_{\QQ_p}$ 的表示。作为 $G(\QQ_p)\times W_{\QQ_p}$ 的表示，我们于是有
\[
H^\bullet_c(\Sht(G,b,\mu)_{\mathbb{C}_p^{\flat}},IC)\otimes_{J_b(\QQ_p)}\pi_\rho\simeq\bigoplus_{\rho'\in(S_{\varphi}/Z(\widehat{G})^{\Gamma})^{\wedge}}\pi_{\rho'}\otimes\sigma_{\rho,\rho',\mu,\varphi},
\]
其中 $IC$ 是 $\Sht(G,b,\mu)$ 的相交上同调复形。当 $\mu$ 极小（minuscule）时，这化为
\[
H^i_c(\Sht(G,b,\mu)_{\mathbb{C}_p^{\flat}},\ol{\QQ}_\ell)\otimes_{J_b(\QQ_p)}\pi_\rho=
\begin{cases}
0 & \text{若 } i\neq d,\\[3pt]
\displaystyle\bigoplus_{\rho'\in(S_{\varphi}/Z(\widehat{G})^{\Gamma})^{\wedge}}\pi_{\rho'}\otimes\sigma_{\rho,\rho',\mu,\varphi}(-d/2) & \text{若 } i=d,
\end{cases}
\]
其中 $d=\dim\,\Sht(G,b,\mu)$。
\end{proposition}

\begin{remark}\label{rem:8.3}
在前述命题中我们用到：$S_{\varphi}$ 的表示 $\mathscr{H}om(\rho,r_\mu\circ\varphi)$ 在 $Z(\widehat{G}^{\Gamma})$ 上是平凡的，因为 $[b]\in B(G,\mu)$ 蕴含 $r_\mu$ 的中心特征由 $\kappa(b)=\widehat{\mu}|_{Z(\widehat{G})^{\Gamma}}$ 给出。
\end{remark}

\srcpage{41}

\subsubsection*{8.2.2 与 $p$-可除群模空间的联系}

最后我们来解释 $\Sht(G,b,\mu)$ 与 Rapoport--Zink 空间的联系。我们把自己放在 §7 那样的 PEL 情形。固定 $\QQ_p$ 上一个局部非分歧 PEL 型的 Rapoport--Zink 数据（[23]）。设 $G$ 是 $\QQ_p$ 上相应的约化群，$\ul{G}$ 是它的约化整模型。\yizhuw{$\ul{G}$ 与 $G$ 的区别（整模型 vs.\ 一般纤维）在本小节反复用到：level 取在 $\ul{G}(\ZZ_p)$ 中，而 $p$-可除群上的 $\ul{G}$-结构记录的是整结构。}固定某个 $b\in G(L)$，对应我们所形变的 $p$-可除群的共变等晶体。此外有一个 Hodge 余特征 $\mu$（与上一小节一样，取它的 Galois 轨道，得到 Rapoport--Zink 空间的非交并）。记
\[
\widehat{\mathcal{M}}
\]
为相应的 Rapoport--Zink 空间，看作 $\Spf(\breve{\ZZ}_p)$ 上的形式概形。记
\[
(\mathcal{M}_K)_{K\subset\ul{G}(\ZZ_p)}
\]
为相应的刚性解析空间塔，$\mathcal{M}_{\ul{G}(\ZZ_p)}=\widehat{\mathcal{M}}_\eta$。于是有 Scholze--Weinstein 的如下定理（[30]）。

\begin{theorem}\label{thm:8.4}
pro-étale 层 $\varprojlim_K\mathcal{M}_K^{\diamond}$ 可由一个完满空间 $\mathcal{M}_\infty$ 表示。此外，$\mathcal{M}_\infty$ 表示大解析 site 上如下预层所关联的层：对 $\breve{\QQ}_p$ 上的 $(R,R^+)$，该预层给出一个元素 $(H,\rho)\in\widehat{\mathcal{M}}(R^+)$，连同 étale $p$-可除群及其 $\ul{G}$-结构 $H\otimes_{R^+}R$ 的一个平凡化。
\end{theorem}

这个陈述里，$(H,\rho)$ 这一对指的是 $p$-可除群连同其模 $\varpi_R$ 约化的刚性化；
\[
\rho:\mathbb{H}\otimes_{\ol{\FF}_q}R^+/\varpi_R\longrightarrow H\otimes_{R^+}R^+/\varpi_R
\]
是与所有 $\ul{G}$-结构相容的拟同源。

\begin{remark}\label{rem:8.5}
上文中，由于 Rapoport--Zink 空间是部分 proper 的，可以把 $R^+$ 换成 $R^{\circ}$。
\end{remark}

\begin{remark}\label{rem:8.6}
Scholze--Weinstein 定理更强，它断言 $\mathcal{M}_\infty\sim\varprojlim_K\mathcal{M}_K$，从而塔中某些函数代数在 $\mathcal{M}_\infty$ 上某些函数代数中稠密。
\end{remark}

我们用到 §6 的结果。有态射
\[
g:\mathcal{M}_\infty\longrightarrow\BT\xrightarrow{u}\mathit{Hecke}^{\mu}
\]
它由前述 $\mathcal{M}_\infty$ 的模描述给出。这里我们再一次修改了 $\BT$ 的定义 6.8，把 $\ul{G}$-结构纳入进来。\yizhuw{$\BT$ 是 $p$-可除群（Barsotti--Tate 群）的叠；按定义 6.8 修改后它同时记录 $\ul{G}$-结构，$g$ 才能落到 $\mathit{Hecke}^{\mu}$ 上。}现在容易核对：利用推论 6.10，$g$ 经由定义 $\Sht(G,b,\mu)$ 的纤维积而分解：
\begin{itemize}
\item $\overleftarrow{h}\circ g$ 经过 $x_b:\Spa(\ol{\FF}_p)\to\Bun_{G,\ol{\FF}_p}$ 分解，这靠的是刚性化 $\rho$：它把 $R^+/\varpi_R$ 上 $p$-可除群的 $F$-等晶体所对应的向量丛与 $\mathscr{E}_b$ 等同起来；
\item $\overrightarrow{h}\circ g$ 经过 $x_1$ 分解，这靠的是 $p$-可除群一般纤维的平凡化。
\end{itemize}
下面这件事现在成了 Scholze--Weinstein 结果的重新解释。

\begin{theorem}\label{thm:8.7}
态射 $\mathcal{M}_\infty\xrightarrow{g}\Sht(G,b,\mu)$ 是同构。
\end{theorem}

\begin{remark}\label{rem:8.8}
这个定理并不给出 $\breve{\QQ}_p$ 上 affinoid 完满的 $(R,R^+)$ 上 $p$-可除群的分类，因为要经过“预层所关联的层”。尽管如此，它给出几何点上的分类，即当 $C|\breve{\QQ}_p$ 代数闭时的 $\oo_C$。话虽如此，这个分类正是 Scholze 与 Weinstein 在 [30] 中为证明其结果而首先证明的东西：$g$ 在几何点上是双射。
\end{remark}
